\subsection{An effective-summand obstruction in the actual iterated three-strand family}

Here the strand rank remains \(r=3\). The integer \(k\) counts
ordinary tensor factors in Blasiak's iteration; \(k=2\) is the
binary algebra used in \fid{F56}/\fid{F61}. This finding tests
compatibility with that actual iteration, rather than an
increasing-strand-rank theorem for the original binary family.

Put
\[
 \delta=v+v^{-1},\qquad D=\delta^2,\qquad
 \Lambda=\mathbb Z[v,v^{-1}],\qquad K=\mathbb Q(v).
\]
We rename the source parameter \(u\) as \(v\).
Definition~2.5 of \cite{BlasiakChebyshev}, p.~6, defines
\[
 \check H_3^{(k)}\subset H_3^{\otimes k},\qquad
 P_i^{(k)}=\delta^k\Delta_i^{(k)}(p_+),\qquad
 Q_i^{(k)}=\delta^k-P_i^{(k)}.
\]
The Hopf structure used here is on the parabolic \(H_2\)
\cite[Proposition~2.3, pp.~4--5]{BlasiakChebyshev}.
The actual restriction inclusions are
\begin{equation}\label{eq:iteratedActualCoproduct}
 \begin{split}
 \check\Delta^{(k_l,k_r)}:\check H_3^{(k_l+k_r)}
     &\hookrightarrow
       \check H_3^{(k_l)}\otimes\check H_3^{(k_r)},\\
 P_i^{(k_l+k_r)}
     &\longmapsto P_i^{(k_l)}\otimes P_i^{(k_r)}
                   +Q_i^{(k_l)}\otimes Q_i^{(k_r)}.
 \end{split}
\end{equation}
These are source equations~(8)--(11), p.~6. No Hopf structure
on the full three-strand algebra is assumed.

In the source's left-column convention,
\(X^{(k)}(c)\) has matrices
\begin{equation}\label{eq:iteratedSourceSimple}
 P_1^{(k)}=\begin{pmatrix}\delta^k&c\\0&0\end{pmatrix},
 \qquad
 P_2^{(k)}=\begin{pmatrix}0&0\\1&\delta^k\end{pmatrix}.
\end{equation}
It is absolutely irreducible when \(c\ne0,\delta^{2k}\)
\cite[\S3.1, equations~(13)--(14), pp.~8--9]{BlasiakChebyshev}.
Transposition converts this to the earlier right-row convention.
Write
\[
 A_j=\delta^jT_j(1/\delta),\qquad
 c_{k,j}=\delta^{2(k-j)}A_j^2,\qquad 1\le j\le k.
\]
Theorem~3.4, p.~12, lists the generic simples as the two
scalar representations and \(X^{(k)}(c_{k,j})\).
Corollary~3.5, p.~13, gives generic split semisimplicity.
Theorem~3.7, pp.~13--14, equations~(27)--(28), gives
\[
 (P_i^{(k)})^2=\delta^kP_i^{(k)},\qquad
 P_1^{(k)}F_k(P_2^{(k)}P_1^{(k)})
 =P_2^{(k)}F_k(P_1^{(k)}P_2^{(k)}),
 \quad F_k(y)=\prod_{j=1}^k(y-c_{k,j}).
\]
These are known source results. The source itself prints
\begin{equation}\label{eq:iteratedChebyshevParameters}
 A_1=1,\quad A_2=2-D,\quad A_3=4-3D,\quad
 A_4=D^2-8D+8
\end{equation}
in equation~(24), p.~11.

\begin{finding}[\fid{F63}, the actual iterated effective-summand boundary]
\label{f:iteratedChebyshevGate}
For every \(k\ge4\), the actual generic simple
\(X^{(k)}(c_{k,4})\) has no basis in which both named
generators have entries in \(\mathbb Q_{\ge0}[v,v^{-1}]\).
Every individual generic simple for \(k\le3\) does have
a model with entries in \(\mathbb N[v,v^{-1}]\).

The \(k=4\) obstruction occurs as a summand, regular at \(v=1\),
in the actual restriction of \(N(D-2)\otimes N(D-2)\)
through \eqref{eq:iteratedActualCoproduct}.
There is an explicit regular commuting projector onto that
summand. No basis makes this projector and both named generators
simultaneously nonnegative Laurent matrices.

Nevertheless every unmarked word trace of the whole
four-dimensional tensor module is a nonnegative integral Laurent
polynomial. Those traces alone therefore do not exclude a
strict effective model for the whole module.
\end{finding}

\begin{proof}
Multiplying \eqref{eq:iteratedSourceSimple} gives
\[
 \operatorname{Tr}(P_1^{(k)}P_2^{(k)})=c.
\]
At \(k=j=4\), put \(A=A_4\). Its invariant two-word trace is
\begin{equation}\label{eq:iteratedNegativeTrace}
 \begin{split}
 A^2={}&v^8-8v^6+12v^4+8v^2+38\\
       &+8v^{-2}+12v^{-4}-8v^{-6}+v^{-8}.
 \end{split}
\end{equation}
A product of nonnegative Laurent matrices has nonnegative
Laurent trace, and simultaneous basis change preserves trace.
The negative coefficients therefore exclude every such basis
of this actual simple.

For all \(k\ge4\), choose \(v_0>1\) with
\[
 D(v_0)=4+2\sqrt2.
\]
Such a value exists since \(D(v)=v^2+2+v^{-2}\)
increases continuously from \(4\) to infinity on \(v>1\).
The right side is a root of \(D^2-8D+8\). Consequently
\[
 c_{k,4}=\delta^{2(k-4)}A^2\ne0,\qquad c_{k,4}(v_0)=0.
\]
A nonzero nonnegative Laurent polynomial is positive at
every positive argument, so \(c_{k,4}\) is not in that cone.
Since \(c_{k,4}(1)=4^{k-1}>0\), its coefficients have both signs.

For \(k\le3\), use only \(j\le k\) among
\[
 b_{k,1}=\delta^{k-1},\qquad
 b_{k,2}=\delta^{k-2}(D-2),\qquad
 b_{k,3}=\delta^{k-3}(3D-4).
\]
These are nonnegative integral Laurent polynomials and
\(b_{k,j}^2=c_{k,j}\). In particular
\(D-2=v^2+v^{-2}\) and \(3D-4=3v^2+2+3v^{-2}\).
Conjugating \eqref{eq:iteratedSourceSimple} by
\(\operatorname{diag}(b_{k,j},1)\) gives
\[
 \begin{pmatrix}\delta^k&b_{k,j}\\0&0\end{pmatrix},
 \qquad
 \begin{pmatrix}0&0\\b_{k,j}&\delta^k\end{pmatrix}.
\]
The scalar simples have generator values \(\delta^k\) and \(0\).
This proves the limited minimality claim about individual
generic simples. It does not supply compatible effective
coproduct intertwiners or regular simultaneous splittings
of all tensor lattices.

For the regular occurrence, use the left-column transpose
of the \fid{F56} \(N(D-2)\) matrices and put
\[
 a=D-2,\qquad d=D^2,\qquad A=D^2-8D+8,
\]
\[
 P_1=\begin{pmatrix}D&a\\0&0\end{pmatrix},\quad
 P_2=\begin{pmatrix}0&0\\a&D\end{pmatrix},\quad
 Q_i=DI_2-P_i.
\]
The actual source restriction gives
\begin{equation}\label{eq:iteratedFourDimensionalMatrices}
 \begin{split}
 R_1&=P_1\otimes P_1+Q_1\otimes Q_1
  =\begin{pmatrix}
     d&Da&Da&2a^2\\
     0&0&0&-Da\\
     0&0&0&-Da\\
     0&0&0&d
    \end{pmatrix},\\
 R_2&=P_2\otimes P_2+Q_2\otimes Q_2
  =\begin{pmatrix}
     d&0&0&0\\
     -Da&0&0&0\\
     -Da&0&0&0\\
     2a^2&Da&Da&d
    \end{pmatrix}.
 \end{split}
\end{equation}
Let \(B\) have columns
\[
 u_+=(d,-Da,-Da,d)^{\mathsf T},\quad
 u_-=(0,1,-1,0)^{\mathsf T},
\]
\[
 b_1=(A,-Da,-Da,d)^{\mathsf T},\quad
 b_2=(d,-Da,-Da,A)^{\mathsf T}.
\]
Their entire action table is
\[
\begin{array}{c|rrrr}
 &u_+&u_-&b_1&b_2\\ \hline
 R_1&d u_+&0&d b_1&A b_1\\
 R_2&d u_+&0&A b_2&d b_2 .
\end{array}
\]
For example, the first coordinate of \(R_1b_2\) is
\(d^2-2da^2+2a^2A=A^2\), since \(A=2a^2-d\);
the other coordinates agree with \(A b_1\).
The remaining products use the same identity.

Subtract \(u_+\) from columns \(b_1,b_2\).
They become \((A-d)e_1\) and \((A-d)e_4\).
The middle two coordinates of \(u_+,u_-\) have determinant
\(2Da\). Since \(A-d=-8(D-1)\),
\begin{equation}\label{eq:iteratedRegularDeterminant}
             \det B=2Da(A-d)^2
                    =128D(D-2)(D-1)^2.
\end{equation}
This is a unit in \(R=\mathbb Q[v,v^{-1}]_{(v-1)}\).
We have therefore proved the actual regular decomposition
\begin{equation}\label{eq:iteratedRegularDecomposition}
 N(D-2)\otimes N(D-2)
       \simeq\check\epsilon_+\oplus\check\epsilon_-\oplus N(A),
\end{equation}
where the plus scalar has value \(d\), the minus scalar
has value zero, and
\[
 N_1(A)=\begin{pmatrix}d&A\\0&0\end{pmatrix},\qquad
 N_2(A)=\begin{pmatrix}0&0\\A&d\end{pmatrix}.
\]
The trace is \(A^2\), which is neither \(0\) nor \(d^2\)
over \(K\). This is the actual \(M^{(4)}\) of the source
classification, with the named matrices and local lattice retained.

Define on the four-dimensional tensor
\begin{equation}\label{eq:iteratedMarkedProjector}
 E_4=\frac{d(R_1+R_2)-R_1R_2-R_2R_1}{d^2-A^2},
 \qquad d^2-A^2=16(D-1)(D-2)^2 .
\end{equation}
Its denominator is a unit in \(R\). The numerator is zero
on both scalar blocks of
\eqref{eq:iteratedRegularDecomposition}, and equals
\((d^2-A^2)I_2\) on \(N(A)\), by direct multiplication.
Thus
\[
 E_4^2=E_4,\quad [E_4,R_i]=0,\quad
 \operatorname{Tr}E_4=2,\quad
 \operatorname{Tr}(E_4R_1R_2)=A^2.
\]
In the original tensor basis it is
\[
 E_4=\frac1{8(D-1)}
 \begin{pmatrix}
 -A&-Da&-Da&-d\\
 Da&d&d&Da\\
 Da&d&d&Da\\
 -d&-Da&-Da&-A
 \end{pmatrix}.
\]
If \(E_4,R_1,R_2\) all had nonnegative Laurent matrices in
one basis, their product trace would contradict
\eqref{eq:iteratedNegativeTrace}. This projector is regular
at \(v=1\).

Finally, let \(w\) be a nonempty word of length \(L\) in
the two generators, and let \(2s\) be its number of cyclic
label changes. Then \(L\ge2s\). Factor
\(N_i=u_i r_i\), where
\[
 u_1=(d,0)^{\mathsf T},\quad r_1=(1,A/d),\qquad
 u_2=(0,d)^{\mathsf T},\quad r_2=(A/d,1).
\]
The scalar \(r_i u_j\) is \(d\) for equal labels and
\(A\) for different labels. The trace closes the final
scalar product back to the initial label. Therefore
\[
 \operatorname{Tr}_{N(A)}(w)=d^{L-2s}A^{2s},
\]
and the two scalar blocks give
\begin{equation}\label{eq:iteratedAllWordTrace}
 \operatorname{Tr}_{N(D-2)^{\otimes2}}(w)
          =d^{L-2s}\bigl(d^{2s}+A^{2s}\bigr).
\end{equation}
Both
\[
 U=d+A=2(D-2)^2,\qquad V=d-A=8(D-1)
\]
are nonnegative integral Laurent polynomials. For \(s\ge1\),
\[
 d^{2s}+A^{2s}
   =2^{1-2s}
      \sum_{\substack{0\le j\le2s\\ j\ \mathrm{even}}}
        \binom{2s}{j}U^{2s-j}V^j.
\]
The right side has nonnegative rational coefficients,
and the left side has integral coefficients. Hence those
coefficients are nonnegative integers. For \(s=0\) the
factor is \(2\). Multiplication by \(d^{L-2s}\) proves
positivity for every nonempty word; the empty trace is \(4\).
\end{proof}

\paragraph{Actual tensor occurrence and its precise boundary.}
The regular comparison in \fid{F21}/\fid{F43} supplies
\(N(D-2)\) summands in each binary rank-three tensor module
with two two-dimensional ordinary factors.
Tensoring one such summand from each of two binary
modules and restricting through
\eqref{eq:iteratedActualCoproduct} gives exactly
\eqref{eq:iteratedFourDimensionalMatrices}.
Thus the tensor cube of four two-dimensional ordinary
factors contains the displayed regular summand after
the stated localization.
Alternatively, the ordinary Specht module \(M_{(2,1)}\)
occurs in each factor's tensor cube, and source
equations~(25)--(26), p.~12, place \(M^{(4)}\) in its
fourth tensor power \cite{BlasiakChebyshev}.

For \(k>4\), append ordinary trivial \(M_{(3)}\) factors.
On such a factor \(P_i^{(1)}=\delta\) and \(Q_i^{(1)}=0\);
the actual restriction map scales the previous generator
by \(\delta\). Its two-word trace is therefore
\(\delta^{2(k-4)}A^2\).
No claim that \(M_{(2,1)}^{\otimes k}\) contains this
label when \(k\) is odd is needed.

A regular linear idempotent does not automatically become
an effective coefficient-kernel operator. Requiring the
projector \(E_4\) to remain effective is precisely the
extra condition excluded above. The unmarked two-word
trace of the whole tensor is, for example,
\[
 d^2+A^2
 =2v^8+40v^4+64v^2+108+64v^{-2}+40v^{-4}+2v^{-8}.
\]
Its positivity and \eqref{eq:iteratedAllWordTrace} do not
prove existence of a simultaneous positive basis or a
coherent iteration.

\paragraph{Pure point traces.}
All exponents of \(c_{k,4}\) are even and its coefficients
have both signs. Under \(v=q^{n/2}\), its roots are
\(q^j\), of even weights \(2j\). In a point complex
pure of weight \(w\), such a root can occur only in
degree \(2j-w\), by BBD~5.1.8, pp.~126--127.
All visible Euler coefficients would therefore have
the same sign \((-1)^w\). This contradicts the mixed
signs just proved. Under \(v=q^n\) the root weights
are again even, giving the same conclusion.
An overall cohomological parity change or common Tate
twist cannot repair it. Bounded mixed Tate complexes
remain possible.

Consequently \fid{F56}/\fid{F61} cannot extend to a
family supplying strict effective finite graded kernels
on every actual simple of every iterated algebra with
the stated word-trace comparison: the \(k=4\) simple
already excludes this. Whole tensor objects whose simple
projectors are not effective are not excluded by this
argument. An independently inserted linear endomorphism
trace is different; the marked pure-trace version
requires the extra comparison turning that effective
closed word into the Frobenius trace object.

\paragraph{What the source's cellular bases assert.}
Definition~5.1 of \cite{BlasiakChebyshev}, p.~20,
requires an anti-involution and triangular multiplication
with coefficients in the specified ring; it imposes no
nonnegative cone.
For \(m=3\), equation~(45), p.~21, uses factors
\(P_{21}-\sigma_j^2\), where
\(\sigma_j=\delta^kT_j(1/\delta)\).
Proposition~5.3, p.~22, requires the indicated
\(\sigma_j\) to be invertible, together with the
earlier differences for incomparable labels.
Equation~(48), p.~22, gives balanced cell matrices
with off-diagonal \(\sigma_j\).
For \(j=k=4\) this is the already printed mixed
parameter \(A_4\).

That localization may omit its positive zero, but it
cannot create a nonnegative Laurent matrix basis:
the generic trace \eqref{eq:iteratedNegativeTrace}
is invariant. There is no contradiction with cellularity,
Propositions~5.5 and~5.8, pp.~23--24, or generic
semisimplicity. No IV/VIII conjecture for these
iterated \(k\ge4\) bases follows from the term
\emph{cellular}. Their signed or factored conventions,
designated bases, and original binary all-rank problem
remain distinct.

% Separate paragraph for Part III.
\paragraph{The iterated obstruction in Grothendieck's language.}
The actual restriction map, regular basis change and
projector are finite locally free module/action data.
EGA~I, Theorem~1.3.7, p.~86, Corollary~1.3.8, p.~87,
and Theorem~1.4.1, pp.~90--92 give the affine module
dictionary \cite{EGA}.
SGA~3, I, 4.7.2--3 and cote~90, scan~17 supply the
action language, not a new Hopf structure on the whole
three-strand algebra \cite{SGA3,Archive}.

Strict graded coefficients use the effective split-torus
character cone \(\mathbb N[v,v^{-1}]\).
The invariant trace tests this cone on the actual module.
SGA~6, V, \S\S2--3 and cote~133, scan~18 provide
representation and \(\lambda/\sigma\) vocabulary for
the symmetric coefficient category; no \(\lambda\)-ring
structure is asserted on the noncommutative word algebra
\cite{SGA6,Archive}.
The pure-point consequence needs the stated
grading-to-Frobenius comparison and BBD~5.1.5/5.1.8,
pp.~125--127 \cite{BBD}.
Mixed bounded complexes remain in SGA~6, I,
Definition~4.7, pp.~121--122 and Example~4.8,
p.~122 \cite{SGA6}.
These framework identifications do not make a regular
linear projector an effective categorical projector;
that is exactly the additional comparison isolated by
\eqref{eq:iteratedMarkedProjector}.

